% ----------
% Tech report template
% 
% ----------
\documentclass[10pt,a4paper]{article}
\usepackage{geometry}
%\usepackage[utf8x]{inputenc}
\usepackage[T1]{fontenc}
\usepackage[english]{babel}
\usepackage[runin]{abstract}
\usepackage{titling}
\usepackage{listings}
\usepackage{booktabs}
\usepackage{mlmodern}
\usepackage{enumitem}
%\usepackage{piton} % Without this enable, it might not work for the listing package. Switch to LuaLateX if you want to use this option, however. 
\usepackage{fancyhdr}
%\usepackage{helvet}
\usepackage{csquotes}
\usepackage{graphicx}
\usepackage{blindtext}
\usepackage{parskip}
\usepackage{etoolbox}
%\usepackage[scaled=0.90]{helvet} % Font 1
%\usepackage[scaled=0.85]{beramono} % Font 2
\usepackage{xcolor}
\usepackage{hyperref}
% Bibliography
\usepackage[backend=biber,style=alphabetic]{biblatex}
\addbibresource{bibb.bib} %Imports bibliography file

% preamble.tex

%\geometry{letterpaper, top=2.2in, headheight=1.7in}
\date{}

\makeatletter

\def\fulltitle#1{\def\@fulltitle{#1}}
\def\runningtitle#1{\def\@runningtitle{#1}}
\def\runningauthor#1{\def\@runningauthor{#1}}
\def\affiliation#1{\def\@affiliation{#1}}
\def\department#1{\def\@department{#1}}
\def\memoid#1{\def\@memoid{#1}}
\def\theyear#1{\def\@theyear{#1}}
\def\mydate#1{\def\@mydate{#1}}

\runningtitle{How to write an effective report} % Short title
\author{Fujimiya Amane} % Full list of authors
\runningauthor{Your name et al.} % Short list of authors
\affiliation{Your University} % Affiliation e.g. University or Company
\department{Your Department} % Department or Office
\memoid{Research Memo 001} % ID of the tech report
\theyear{2024} % year of the tech report
\mydate{October, 2024} %the date

\def\displaymydate{\@mydate}
\def\displaytheyear{\@theyear}
\def\displaymemoid{\@memoid}
\def\displaydepartment{\@department}
\def\displayaffiliation{\@affiliation}
\def\displayrunningauthor{\@runningauthor}
\def\displayrunningtitle{\@runningtitle}

\makeatother

\makeatletter
\patchcmd{\@zfancyhead}{\fancy@reset}{\f@nch@reset}{}{}
\patchcmd{\@set@em@up}{\f@ncyolh}{\f@nch@olh}{}{}
\patchcmd{\@set@em@up}{\f@ncyolh}{\f@nch@olh}{}{}
\patchcmd{\@set@em@up}{\f@ncyorh}{\f@nch@orh}{}{}
\makeatother

% ----------
% fancyhdr setup
% ----------

\pagestyle{fancy}
\renewcommand{\headrulewidth}{0pt}

\fancypagestyle{firstpage}
{
    \fancyhf{}
    \fancyhead{%
        \begin{center}
            \textsc{\displayaffiliation}\\
            \textsc{\displaydepartment}
        \end{center}
        \vspace*{0.5in}
        \displaymemoid
        \hfill
        \displaymydate
    }
    \fancyfoot[L]{\footnotesize \textbf{\textsc{\displayaffiliation} -- \displaytheyear}}
}
\fancyhead[EL]{\small \displayrunningauthor}
\fancyhead[ER]{\small Page \thepage}
\fancyhead[OL]{\small Page \thepage}
\fancyhead[OR]{\small \displayrunningtitle}

\chead{}
\lfoot{}
\cfoot{}
\rfoot{}

% ----------

\providecommand{\keywords}[1]{\noindent \textbf{Keywords:} #1} % keywords definition

%\renewcommand{\familydefault}{\sfdefault} % sans-serif font
\setlength{\parskip}{1em} % 1em skip between paragraphs
\abslabeldelim{:} % use colon after the run-in 'Abstract' word
\setlength{\parindent}{0pt} % no indentantion for first page 
\setlength{\abstitleskip}{-\absparindent} % this is because the abstract package sets \absparindent to \parindent

\pretitle{\Large \begin{center}}
\posttitle{\\[2ex] \Large \textbf{by}\end{center}}

% ----------
% Variables
% ----------

\title{\textbf{Research abstract},\\ v1, first edition} % Full title of your tech report
\runningtitle{Research Record Studies} % Short title
\author{Fujimiya Amane} % Full list of authors
\runningauthor{Amane et al.} % Short list of authors
\affiliation{} % Affiliation e.g. University or Company
\department{Your Department} % Department or Office
\memoid{RES-RP-ABS-XX} % ID of the tech report
\theyear{Year} % year of the tech report
\mydate{Full month, year} %the date
\usepackage{tcolorbox}
\usepackage{fancyvrb}

\tcbuselibrary{minted,breakable,xparse,skins}

\definecolor{bg}{gray}{0.95}
%\DeclareTCBListing{mintedbox}{O{}m!O{}}{%
%  breakable=true,
%  listing engine=minted,
%  listing only,
%  minted language=#2,
%  minted style=default,
%  minted options={%
%    linenos,
%    gobble=0,
%    breaklines=true,
%    breakafter=,,
%    fontsize=\small,
%    numbersep=8pt,
%    #1},
%  boxsep=0pt,
%  left skip=0pt,
%  right skip=0pt,
%  left=25pt,
%  right=0pt,
%  top=3pt,
%  bottom=3pt,
%  arc=5pt,
%  leftrule=0pt,
%  rightrule=0pt,
%  bottomrule=2pt,
%  toprule=2pt,
%  colback=bg,
%  colframe=orange!70,
%  enhanced,
%  overlay={%
%    \begin{tcbclipinterior}
%    \fill[orange!20!white] (frame.south west) rectangle ([xshift=20pt]frame.north west);
%    \end{tcbclipinterior}},
%  #3}

% Python preset environment. 
\newcommand{\pyobject}[1]{\fbox{\color{red}{\texttt{#1}}}}

\NewDocumentCommand{\codeword}{v}{%
\texttt{\textcolor{blue}{#1}}%
}
% Inline python highlighting (Might not look great, but eh)
\lstset{language=Python,keywordstyle={\bfseries \color{blue}}}


% ----------
% actual document
% ----------
\begin{document}
\maketitle

\begin{abstract}
    \noindent
    This is a template for all RHINELAB-RINALAB \textbf{research abstract} provided for \textit{preliminary drafting} and \textit{short-form idea crafting} only (for reference, it will look like those Call for Papers that ask for abstract). 
    
    This is \textbf{version 1} of the template. Rather would there not be a lot of changes, but just so happens, then here. 

    \keywords{Keyword1, Keyword2, Keyword3}
\end{abstract}

\vspace{2.5cm}



\thispagestyle{firstpage}

\pagebreak

% ----------
% End of first page
% ----------

\newgeometry{} % Redefine geometries (normal margins)

\section{Specification}
\label{sec:spec}

Provide specification details, metaprogramming. 

\subsection{Document code, priority}

Provide document codes, additional document code, purpose, organization, authorship and anyone contributing to the report, identification of purpose (proposal, or research, indexing 01). Priority code, additional conditions. Does it override anything, or add to anything?

\subsection{Body of issue}

Who issue, team or departmental attachment. 

\subsection{Date and Version}

Version specification across versions, and similarly date of construction. \textbf{Do not remove old dates}, for versioning record. 

\section{Introduction}

What is the topic of the short-form research abstract? Expanded name if able\footnote{This is because most of the time the name of the main document up in the title, will be shortened. This is in procession to hold the full title name before truncation.}.

\subsection{Topics}

Layout relevant topics. If the topic is with organization codes (for example, AMS Mathematics field specification code, or Physics catalogue), please specify. Otherwise, proceed with your own categorization. 

\subsection{Additional topic consideration}

What is the aim or short purpose to those fields, and other short considerations? 

\section{Abstract details}

Here is the main abtract

\subsection{Foundational knowledge}

List your preliminary foundational knowledge required or wanted to write. 

\subsection{Target of research}

What is your target for research? Brief discussion and what you are doing toward the target of research. Do you have: 
\begin{itemize}
    \item Any novel idea?
    \item Any attack of consideration against the target?
    \item Notable differences and interjections\footnote{Here, interjection is used to mean \textit{The act of throwing or inserting a comment, remark, or objection into a conversation, debate, or speech on opposition to the current `continuous' idea or framework}. The different from the normal usage of the word here, is that it means treating current treatise and working theories as the \textit{continuous speech} to be interrupted of the person writing this abstract.}?
    \item Comparison and learned or analyzed insights?   
\end{itemize}
Among other things, please write here. 
\subsection{Evaluation}
Is it: 
\begin{itemize}
    \item New? 
    \item Novel? 
    \item Possible? 
    \item Feasible and deployable? 
    \item Impactful? 
    \item Necessity? 
    \item Etc.
\end{itemize}
Then, what is it? Why? How? Write here. 
\section{Conclusion}
Write your own conclusion. 






\clearpage
% ----------
% Bibliography
% ----------

% \bibliography{../biblio.bib}
% \bibliographystyle{abbrv}

\appendix

\section{An appendix}

Appendix to more information. Analysis, cost process, etc. 

\section{Another appendix}

\subsection*{Subsectioning in appendix}

Some more text, and a list:

\blindenumerate



\section{Citations}
Items that are cited: \textit{The \LaTeX\ Companion} book \cite{latexcompanion} together with Einstein's journal paper \cite{einstein} and Dirac's book \cite{dirac}---which are physics-related items. Next, citing two of Knuth's books: \textit{Fundamental Algorithms} \cite{knuth-fa} and \textit{The Art of Computer Programming} \cite{knuthwebsite}.


Remember to change this to your respective bibliography. Usually, the file that holds the bibliography in the path is \texttt{bibb.bib}, so beware of that. The below part where there is \texttt{printbibliography} does it work - it prints the bibliography. Cite them using the \texttt{cite} key seen in the above citation. 

\medskip

\printbibliography %Prints bibliography


\end{document}

% ----------
